\documentclass[12pt,a4paper]{report}
% Impostazioni lingua e codifica
\usepackage[utf8]{inputenc}
\usepackage[T1]{fontenc}
\usepackage[italian]{babel}
% Matematica e simboli
\usepackage{amsmath, amssymb, amsthm}
\usepackage{mathtools}
% Figure e grafica
\usepackage{graphicx}
\usepackage{caption}
\usepackage{subcaption}
% Impostazioni pagina e collegamenti
\usepackage{geometry}
\geometry{margin=1.5cm}
\usepackage{hyperref}
\hypersetup{
colorlinks=true,
linkcolor=blue,
filecolor=magenta,
urlcolor=cyan,
pdfauthor={Dispensa Universitaria - Enrico Desideri},
pdftitle={Biochimica}
}
% Ambienti personalizzati per definizioni e teoremi
\newtheorem{definizione}{Definizione}[section]
\newtheorem{teorema}{Teorema}[section]
\newtheorem{proposizione}{Proposizione}[section]
\newtheorem{esempio}{Esempio}[section]
\newtheorem{nota}{Nota}[section]

\begin{document}

\title{\textbf{Biochimica}}
\maketitle
\tableofcontents
\clearpage

% =========================================================
% CAPITOLO 1: Organizzazione della Materia
% =========================================================
\chapter{Organizzazione della Materia}
\section{La materia e gli elementi chimici}
La \textbf{materia} è definita come qualunque entità che occupi uno spazio e possieda una massa. Essa è costituita da \textbf{elementi chimici} o da combinazioni di essi. Attualmente sono noti 118 elementi, di cui 92 presenti naturalmente sulla Terra; i restanti sono stati sintetizzati artificialmente in laboratorio. Ciascun elemento è identificato da un simbolo composto da una o due lettere (ad esempio, \(\mathrm{O}\) per l'ossigeno, \(\mathrm{Na}\) per il sodio).
Nel corpo umano sono presenti meno di 30 dei 92 elementi naturali, distribuiti in proporzioni molto diverse. Quattro elementi --- ossigeno (\(\mathrm{O}\)), carbonio (\(\mathrm{C}\)), idrogeno (\(\mathrm{H}\)) e azoto (\(\mathrm{N}\)) --- rappresentano oltre il 96\% del peso corporeo. La maggior parte dell'ossigeno è contenuta nell'acqua (\(\mathrm{H_2O}\)), che costituisce il 50--70\% della massa corporea totale.
Altri elementi presenti in quantità significative includono:
\begin{itemize}
    \item Calcio (\(\mathrm{Ca}\)): 1,5\%
    \item Fosforo (\(\mathrm{P}\)): 1\%
    \item Potassio (\(\mathrm{K}\)): 0,35\%
    \item Zolfo (\(\mathrm{S}\)): 0,25\%
    \item Sodio (\(\mathrm{Na}\)): 0,2\%
    \item Cloro (\(\mathrm{Cl}\)): 0,2\%
    \item Ferro (\(\mathrm{Fe}\)): 0,005\%
\end{itemize}
Il restante 0,2\% è costituito da \textbf{elementi in tracce}, principalmente metalli come rame (\(\mathrm{Cu}\)), zinco (\(\mathrm{Zn}\)) e manganese (\(\mathrm{Mn}\)), essenziali per funzioni enzimatiche e strutturali.
Nella tavola periodica, gli elementi sono organizzati in \textbf{gruppi} (colonne) e \textbf{periodi} (righe). Elementi appartenenti allo stesso gruppo condividono proprietà chimiche simili, poiché presentano lo stesso numero di elettroni di valenza.

\section{Struttura dell'atomo}
L'\textbf{atomo} è la più piccola unità di un elemento che ne conserva le proprietà chimiche. Esso può esistere isolato o legarsi ad altri atomi, uguali o diversi, per formare molecole.
Ogni atomo è costituito da:
\begin{itemize}
    \item \textbf{Nucleo centrale}, contenente:
    \begin{itemize}
        \item \textit{Protoni} (\(p^+\)): particelle con carica elettrica positiva;
        \item \textit{Neutroni} (\(n^0\)): particelle neutre, di massa simile ai protoni.
    \end{itemize}
    \item \textbf{Elettroni} (\(e^-\)): particelle con carica negativa, di massa circa 2000 volte inferiore a quella di protoni e neutroni, che si muovono in regioni dello spazio dette \textit{orbitali elettronici}.
\end{itemize}
Gli orbitali sono organizzati in livelli energetici progressivamente più esterni. Il primo orbitale (più vicino al nucleo) può contenere al massimo 2 elettroni, il secondo fino a 8, il terzo fino a 18, e così via.
\begin{definizione}[Numero atomico]
Il \textbf{numero atomico} (\(Z\)) di un elemento è il numero di protoni presenti nel nucleo di ciascun suo atomo. Poiché in un atomo neutro il numero di elettroni eguaglia quello dei protoni, il numero atomico determina anche la configurazione elettronica e, di conseguenza, le proprietà chimiche dell'elemento.
\end{definizione}
Ad esempio:
\begin{itemize}
    \item Idrogeno (\(\mathrm{H}\)): \(Z = 1\) (1 protone, 1 elettrone);
    \item Ossigeno (\(\mathrm{O}\)): \(Z = 8\) (8 protoni, 8 elettroni);
    \item Uranio (\(\mathrm{U}\)): \(Z = 92\) (elemento naturale con il più alto numero atomico).
\end{itemize}

\section{Isotopi e massa atomica}
Tutti gli atomi di uno stesso elemento condividono il numero atomico, ma possono differire per il numero di neutroni nel nucleo. La somma di protoni e neutroni è detta \textbf{numero di massa} (\(A\)).
\begin{definizione}[Isotopi]
Gli \textbf{isotopi} sono atomi dello stesso elemento (stesso \(Z\)) con diverso numero di massa (\(A\)), ovvero con differente numero di neutroni.
\end{definizione}
Un esempio classico è rappresentato dagli isotopi dell'idrogeno:
\begin{itemize}
    \item \textbf{Prozio} (\(^1\mathrm{H}\)): 1 protone, 0 neutroni (isotopo più abbondante);
    \item \textbf{Deuterio} (\(^2\mathrm{H}\)): 1 protone, 1 neutrone (stabile);
    \item \textbf{Trizio} (\(^3\mathrm{H}\)): 1 protone, 2 neutroni (radioisotopo instabile).
\end{itemize}
La maggior parte degli isotopi è stabile; alcuni, detti \textbf{radioisotopi}, sono instabili e decadono emettendo radiazioni. Gli isotopi di un elemento condividono le stesse proprietà chimiche, poiché la configurazione elettronica --- determinante per le reazioni chimiche --- è identica.
Nella tavola periodica, per ciascun elemento è riportata la \textbf{massa atomica relativa}, espressa in \textbf{unità di massa atomica} (u.m.a.) o \textbf{dalton} (\(\mathrm{Da}\)). Un dalton è definito come \(1/12\) della massa di un atomo di carbonio-12 (\(^{12}\mathrm{C}\)). La massa atomica indicata è la media ponderata delle masse di tutti gli isotopi naturali, in base alla loro abbondanza relativa.

\section{Ioni e ionizzazione}
Un atomo neutro può acquisire o perdere uno o più elettroni, trasformandosi in uno \textbf{ione} carico elettricamente.
\begin{itemize}
    \item \textbf{Catione}: ione con carica positiva, formato per perdita di elettroni (es. \(\mathrm{Ca^{2+}}\): calcio che ha perso 2 elettroni);
    \item \textbf{Anione}: ione con carica negativa, formato per acquisizione di elettroni (es. \(\mathrm{Cl^-}\): cloro che ha acquisito 1 elettrone).
\end{itemize}
La ionizzazione è un processo fondamentale in molte reazioni biochimiche, come la formazione di sali e il trasporto ionico attraverso le membrane cellulari.

\section{Legami chimici}
Gli atomi si legano tra loro per formare molecole, raggiungendo configurazioni elettroniche più stabili (a minore energia potenziale). La formazione o rottura di legami chimici comporta sempre scambi di energia.
\subsection{Legame covalente}
Nel \textbf{legame covalente}, due atomi condividono una o più coppie di elettroni. È il legame più forte tra quelli rilevanti in biologia.
\begin{definizione}[Valenza]
La \textbf{valenza} di un atomo è il numero di elettroni che esso può mettere in condivisione per formare legami covalenti.
\end{definizione}
Esempi:
\begin{itemize}
    \item Idrogeno (\(\mathrm{H}\)): valenza 1;
    \item Ossigeno (\(\mathrm{O}\)): valenza 2;
    \item Carbonio (\(\mathrm{C}\)): valenza 4.
\end{itemize}
A seconda del numero di coppie condivise, si distinguono:
\begin{itemize}
    \item \textbf{Legame singolo}: condivisione di 1 coppia di elettroni;
    \item \textbf{Legame doppio}: condivisione di 2 coppie;
    \item \textbf{Legame triplo}: condivisione di 3 coppie.
\end{itemize}
Maggiore è il numero di elettroni condivisi, maggiore è l'\textbf{energia di legame} (espressa in \(\mathrm{kJ\,mol^{-1}}\) o \(\mathrm{kcal\,mol^{-1}}\)), ovvero l'energia necessaria per rompere il legame.
\subsubsection{Polarità del legame covalente}
La distribuzione degli elettroni condivisi dipende dall'\textbf{elettronegatività}, ovvero la capacità di un atomo di attrarre elettroni.
\begin{itemize}
    \item \textbf{Legame covalente non polare}: gli atomi hanno elettronegatività simile; gli elettroni sono equamente condivisi (es. \(\mathrm{O_2}\), \(\mathrm{CH_4}\)).
    \item \textbf{Legame covalente polare}: differenza di elettronegatività \(0 < \Delta e < 1.9\); l'atomo più elettronegativo acquisisce una parziale carica negativa (\(\delta^-\)), l'altro una parziale carica positiva (\(\delta^+\)). Esempio: legame \(\mathrm{O-H}\) nell'acqua.
\end{itemize}
\begin{definizione}[Molecola polare]
Una molecola è \textbf{polare} quando presenta una separazione permanente di carica dovuta a legami polari e a una geometria asimmetrica. Al contrario, una molecola \textbf{apolare} ha distribuzione simmetrica degli elettroni.
\end{definizione}
La polarità influenza la solubilità: le molecole polari sono \textbf{idrofiliche} (si sciolgono in acqua), mentre quelle apolari sono \textbf{idrofobiche} (es. lipidi, oli).
\subsection{Legame ionico}
Quando la differenza di elettronegatività tra due atomi è elevata (\(\Delta e > 1.9\)), un atomo trasferisce completamente uno o più elettroni all'altro. Si formano ioni di carica opposta che si attraggono elettrostaticamente, dando origine al \textbf{legame ionico}.
Esempio: nel cloruro di sodio (\(\mathrm{NaCl}\)), il sodio cede un elettrone al cloro, formando \(\mathrm{Na^+}\) e \(\mathrm{Cl^-}\) che si organizzano in un \textit{reticolo cristallino}. La carica complessiva del composto è sempre nulla.
\subsection{Legame idrogeno}
Il \textbf{legame idrogeno} (o ponte a idrogeno) è un'interazione debole (10--20 volte meno energetica di un legame covalente) che si instaura tra molecole polari contenenti idrogeno legato a un atomo altamente elettronegativo (ossigeno, azoto, fluoro).
Si forma quando la parziale carica positiva dell'idrogeno (\(\delta^+\)) interagisce con la parziale carica negativa (\(\delta^-\)) di un'altra molecola. Nell'acqua, ogni molecola può formare fino a 4 legami idrogeno, creando una rete dinamica responsabile di proprietà uniche come coesione, adesione e alto calore specifico.
Nonostante la debolezza individuale, la grande quantità di legami idrogeno nelle macromolecole biologiche (proteine, DNA) conferisce stabilità strutturale e, al contempo, flessibilità dinamica, essenziale per la funzione biologica.

\section{Molecole, gruppi funzionali e composti organici}
Le molecole sono aggregati di atomi legati covalentemente. Quando contengono atomi di due o più elementi diversi, sono dette \textbf{composti}.
\subsection{Gruppi funzionali}
I \textbf{gruppi funzionali} sono specifici raggruppamenti di atomi che conferiscono caratteristiche chimiche definite alle molecole in cui sono presenti. Nei composti biologici, i più comuni sono:
\begin{itemize}
    \item \textbf{Carbossilico} (\(-\mathrm{COOH}\)): acido, presente negli amminoacidi;
    \item \textbf{Aminico} (\(-\mathrm{NH_2}\)): basico, presente negli amminoacidi;
    \item \textbf{Idrossilico} (\(-\mathrm{OH}\)): polare, presente in zuccheri e alcoli;
    \item \textbf{Fosfato} (\(-\mathrm{PO_4^{2-}}\)): coinvolto nel trasferimento di energia (ATP) e nella fosforilazione;
    \item \textbf{Carbonilico} (\(>\mathrm{C=O}\)): presente in aldeidi e chetoni (es. monosaccaridi);
    \item \textbf{Sulfidrilico} (\(-\mathrm{SH}\)): importante per la struttura terziaria delle proteine (ponti disolfuro nella cisteina).
\end{itemize}
\subsection{Composti organici e scheletro carbonioso}
I \textbf{composti organici} sono definiti come molecole contenenti atomi di carbonio (\(\mathrm{C}\)), con eccezione di ossidi del carbonio e carbonati. Le quattro classi di macromolecole biologiche --- carboidrati, lipidi, proteine e acidi nucleici --- sono tutte composti organici.
Il carbonio, con valenza 4, forma legami covalenti stabili con sé stesso e con altri elementi, dando origine a:
\begin{itemize}
    \item Catene lineari;
    \item Strutture ramificate;
    \item Anelli ciclici.
\end{itemize}
Questa versatilità è alla base della diversità strutturale delle molecole biologiche.
\subsection{Notazioni molecolari}
Per rappresentare le molecole si utilizzano diverse convenzioni:
\begin{itemize}
    \item \textbf{Formula molecolare (bruta)}: indica il tipo e il numero di atomi (es. glucosio: \(\mathrm{C_6H_{12}O_6}\)). Permette il calcolo del peso molecolare:
    \[
    \text{PM}_{\text{glucosio}} = (12 \times 6) + (1 \times 12) + (16 \times 6) = 180\,\mathrm{Da}
    \]
    \item \textbf{Formula di struttura}: mostra come gli atomi sono collegati tra loro, fornendo informazioni sulla connettività e sulla geometria.
\end{itemize}
\begin{definizione}[Isomeri]
Gli \textbf{isomeri} sono molecole con la stessa formula molecolare ma diversa formula di struttura o diversa disposizione spaziale degli atomi.
\end{definizione}
\subsubsection{Isomeria strutturale}
Gli isomeri strutturali differiscono per il numero o il tipo di legami tra gli atomi:
\begin{itemize}
    \item \textbf{Isomeri di catena}: differenze nello scheletro carbonioso (es. leucina e isoleucina);
    \item \textbf{Isomeri di posizione}: diversa posizione di gruppi funzionali o legami multipli (es. acidi grassi insaturi).
\end{itemize}
\subsubsection{Stereoisomeria}
Gli stereoisomeri hanno gli stessi legami ma diversa disposizione tridimensionale:
\begin{itemize}
    \item \textbf{Isomeri geometrici (cis-trans)}: in presenza di doppi legami \(\mathrm{C=C}\), i sostituenti possono trovarsi dallo stesso lato (\textit{cis}) o da lati opposti (\textit{trans}). Il doppio legame impedisce la rotazione, "bloccando" la conformazione.
    \item \textbf{Isomeri ottici (enantiomeri)}: molecole che sono immagini speculari non sovrapponibili. Si verificano quando un atomo di carbonio (detto \textit{centro chirale}) è legato a quattro gruppi diversi. Gli enantiomeri sono indicati come \textit{D} (destrogiro) e \textit{L} (levogiro).
\end{itemize}
\begin{nota}[Chiralità in biologia]
In sistemi biologici, gli enantiomeri non sono intercambiabili. Nelle proteine, tutti gli amminoacidi sono nella forma \textit{L}; negli zuccheri metabolici, prevale la forma \textit{D} (es. \textit{D}-glucosio). Molti enzimi mostrano specificità stereochimica, riconoscendo solo uno dei due enantiomeri.
\end{nota}

\section{Reazioni chimiche ed energia}
Una \textbf{reazione chimica} comporta la rottura di legami nei reagenti e la formazione di nuovi legami nei prodotti. L'insieme di tutte le reazioni chimiche in un organismo costituisce il \textbf{metabolismo}.
\subsection{Conservazione della massa e dell'energia}
\begin{itemize}
    \item \textbf{Legge di conservazione della massa}: il numero totale di atomi di ciascun elemento è identico tra reagenti e prodotti.
    \item \textbf{Legge di conservazione dell'energia}: l'energia non può essere creata né distrutta, ma solo trasformata da una forma all'altra.
\end{itemize}
L'energia immagazzinata nei legami chimici è \textbf{energia potenziale chimica}, che può essere convertita in \textbf{energia cinetica} (movimento) o \textbf{energia termica} (calore).
\subsection{Unità di misura dell'energia}
\begin{itemize}
    \item \textbf{Joule} (\(\mathrm{J}\)): unità SI; \(1\,\mathrm{J} = 1\,\mathrm{N \cdot m}\);
    \item \textbf{Caloria} (\(\mathrm{cal}\)): calore necessario per innalzare di \(1^\circ\mathrm{C}\) la temperatura di \(1\,\mathrm{g}\) di acqua;
    \item Relazione: \(1\,\mathrm{cal} = 4.184\,\mathrm{J}\); \(1\,\mathrm{kcal} = 1000\,\mathrm{cal} = 4.184\,\mathrm{kJ}\).
\end{itemize}
Un adulto consuma circa \(2000\,\mathrm{kcal}\) al giorno, equivalenti a \(8.4\,\mathrm{MJ}\).
\subsection{Reazioni esoergoniche ed endoergoniche}
\begin{itemize}
    \item \textbf{Reazioni esoergoniche}: l'energia dei prodotti è inferiore a quella dei reagenti; rilasciano energia (es. degradazione di glucosio o lipidi).
    \item \textbf{Reazioni endoergoniche}: l'energia dei prodotti è superiore; richiedono apporto energetico (es. sintesi di ATP, biosintesi di macromolecole).
\end{itemize}
L'\textbf{ATP} (adenosina trifosfato) funge da ``moneta energetica'': la sua idrolisi (\(\mathrm{ATP \rightarrow ADP + P_i}\)) rilascia circa \(30\,\mathrm{kJ/mol}\), utilizzati per lavoro cellulare (contrazione muscolare, trasporto attivo, sintesi).
\subsection{Classificazione delle reazioni biochimiche}
\begin{enumerate}
    \item \textbf{Sintesi (anabolismo)}: due molecole si combinano per formarne una più complessa, spesso con rilascio di acqua (\textit{reazione di condensazione}). Esempio: formazione del legame peptidico tra amminoacidi. Reazioni endoergoniche.
    \item \textbf{Degradazione (catabolismo)}: una molecola complessa viene scissa in unità più semplici, spesso con consumo di acqua (\textit{reazione di idrolisi}). Esempio: idrolisi dell'ATP. Reazioni esoergoniche.
    \item \textbf{Fosforilazione/defosforilazione}: trasferimento di un gruppo fosfato (\(\mathrm{PO_4^{2-}}\)) da o verso una molecola. Meccanismo chiave nella regolazione enzimatica e nell'attivazione metabolica (es. fosforilazione del glucosio nella glicolisi).
    \item \textbf{Ossidoriduzioni (redox)}: trasferimento di elettroni da una molecola (che si \textit{ossida}) a un'altra (che si \textit{riduce}). Spesso coinvolgono il trasferimento di atomi di idrogeno (\textit{deidrogenazioni}). Sono fondamentali nella respirazione cellulare e nella fotosintesi.
\end{enumerate}
\subsection{Reversibilità delle reazioni}
Le reazioni possono essere:
\begin{itemize}
    \item \textbf{Irreversibili}: procedono in una sola direzione;
    \item \textbf{Reversibili}: indicate con doppia freccia (\(\rightleftharpoons\)); le due direzioni avvengono simultaneamente, ma una può essere favorita in base a concentrazione dei reagenti, pH, temperatura, ecc.
\end{itemize}

\section{L'acqua: proprietà e ruolo biologico}
L'acqua (\(\mathrm{H_2O}\)) è il composto inorganico più abbondante negli organismi viventi, costituendo oltre il 60\% della massa corporea (con variazioni tissutali: sangue 90\%, muscolo 75\%, osso 25\%, tessuto adiposo 5\%). Circa due terzi dell'acqua corporea sono intracellulari.
\subsection{Funzioni biologiche dell'acqua}
\begin{itemize}
    \item Trasporto di nutrienti e ossigeno (tramite il sangue);
    \item Eliminazione di prodotti di scarto (urea, sali) tramite le urine;
    \item Termoregolazione mediante sudorazione;
    \item Solvente per reazioni biochimiche;
    \item Reagente (idrolisi) o prodotto (condensazione) in reazioni metaboliche.
\end{itemize}
Una perdita di acqua pari al 2\% della massa corporea compromette già le performance fisiche e cognitive.
\subsection{Polarità e capacità solvente}
La geometria angolare e la differenza di elettronegatività tra ossigeno e idrogeno conferiscono all'acqua una forte polarità: l'ossigeno presenta parziale carica negativa (\(\delta^-\)), gli idrogeni parziale carica positiva (\(\delta^+\)).
Ciò permette all'acqua di:
\begin{itemize}
    \item Formare fino a 4 legami idrogeno per molecola, creando una rete dinamica;
    \item Sciogliere sostanze polari o ioniche (\textbf{idrofiliche}), come zuccheri e sali;
    \item Escludere molecole apolari (\textbf{idrofobiche}), come lipidi, che tendono ad aggregarsi per minimizzare il contatto con l'acqua (\textit{effetto idrofobico}).
\end{itemize}
\begin{definizione}[Molecole anfipatiche]
Molecole che possiedono sia una regione polare/idrofilica sia una regione apolare/idrofobica sono dette \textbf{anfipatiche} (es. fosfolipidi, acidi biliari). In ambiente acquoso, si organizzano spontaneamente in strutture come micelle o doppi strati lipidici, fondamentali per le membrane biologiche.
\end{definizione}
\subsection{Stati fisici dell'acqua}
I cambiamenti di stato riflettono variazioni nella rete di legami idrogeno:
\begin{itemize}
    \item \textbf{Ghiaccio}: molecole fisse, ciascuna forma 4 legami idrogeno stabili;
    \item \textbf{Acqua liquida}: legami idrogeno in continuo formarsi e rompersi;
    \item \textbf{Vapore}: assenza di legami idrogeno significativi.
\end{itemize}

\section{Acidi, basi e regolazione del pH}
\subsection{Definizioni di acido e base}
\begin{itemize}
    \item \textbf{Acido}: composto che, in soluzione acquosa, dona uno ione idrogeno (\(\mathrm{H^+}\), protone). Esempio: acido lattico \(\rightleftharpoons\) lattato + \(\mathrm{H^+}\).
    \item \textbf{Base}: composto che accetta uno ione \(\mathrm{H^+}\). Esempio: \(\mathrm{NH_3 + H^+ \rightleftharpoons NH_4^+}\).
\end{itemize}
Acidi e basi reagiscono tra loro formando sali e acqua (reazione di neutralizzazione): 
\[
\mathrm{HCl + NaOH \rightarrow NaCl + H_2O}
\]
\subsection{Autoionizzazione dell'acqua}
L'acqua può comportarsi sia da acido che da base, autoionizzandosi:
\[
\mathrm{H_2O + H_2O \rightleftharpoons H_3O^+ + OH^-}
\]
In forma semplificata:
\[
\mathrm{H_2O \rightleftharpoons H^+ + OH^-}
\]
\subsection{Scala del pH}
L'acidità o basicità di una soluzione è quantificata dal \textbf{pH}, definito come:
\[
\mathrm{pH} = -\log_{10}[\mathrm{H^+}]
\]
dove \([\mathrm{H^+}]\) è la concentrazione molare di ioni idrogeno.
\begin{itemize}
    \item \(\mathrm{pH} = 7\): soluzione neutra (\([\mathrm{H^+}] = [\mathrm{OH^-}]\), acqua pura);
    \item \(\mathrm{pH} < 7\): soluzione acida;
    \item \(\mathrm{pH} > 7\): soluzione basica/alcalina.
\end{itemize}
La scala è logaritmica: una soluzione a pH 6 ha una concentrazione di \(\mathrm{H^+}\) dieci volte maggiore di una a pH 7.
\subsection{Importanza fisiologica del pH}
Il pH ematico è strettamente regolato tra 7.35 e 7.45. Valori esterni a 7.0--7.8 possono essere letali. Durante esercizio intenso, la produzione di acido lattico può abbassare il pH muscolare fino a ~6.5, compromettendo la contrattilità.
\subsection{Sistemi tampone}
Per contrastare variazioni di pH, l'organismo utilizza \textbf{sistemi tampone}, costituiti da una coppia acido debole/base coniugata:
\[
\mathrm{HA \rightleftharpoons H^+ + A^-}
\]
\begin{itemize}
    \item In presenza di eccesso di \(\mathrm{H^+}\): la base coniugata (\(\mathrm{A^-}\)) li lega, riformando l'acido debole (\(\mathrm{HA}\));
    \item In presenza di eccesso di \(\mathrm{OH^-}\): l'acido debole (\(\mathrm{HA}\)) cede \(\mathrm{H^+}\), neutralizzando gli ioni idrossido.
\end{itemize}
\textbf{Principali sistemi tampone fisiologici}:
\begin{itemize}
    \item \textbf{Acido carbonico/bicarbonato} (\(\mathrm{H_2CO_3/HCO_3^-}\)): principale tampone extracellulare (sangue);
    \item \textbf{Fosfati} (\(\mathrm{H_2PO_4^-/HPO_4^{2-}}\)): tampone intracellulare e urinario;
    \item \textbf{Carnosina} e \textbf{fosfocreatina}: tamponi intracellulari nel muscolo scheletrico;
    \item \textbf{Emoglobina}: tampona gli \(\mathrm{H^+}\) nei globuli rossi durante il trasporto di \(\mathrm{CO_2}\).
\end{itemize}
Questi sistemi cooperano per mantenere l'omeostasi acido-base anche in condizioni metaboliche estreme.
% =========================================================
% Fine Capitolo 1
% =========================================================

% =========================================================
% CAPITOLO 2: La cellula eucariotica
% =========================================================
\chapter{La cellula eucariotica}
\section{Introduzione alla cellula}
La \textbf{cellula} è la più piccola unità funzionale in grado di mantenere le caratteristiche del tessuto di appartenenza. Gli organismi viventi possono essere classificati in:
\begin{itemize}
    \item \textbf{Unicellulari}: costituiti da una singola cellula (es. batteri, lieviti);
    \item \textbf{Multicellulari}: composti da numerose cellule specializzate (es. organismi animali e vegetali).
\end{itemize}
Il corpo umano è costituito da circa \(3 \times 10^{13}\) cellule, appartenenti a circa 200 tipi diversi, che si organizzano secondo una gerarchia strutturale:
\[
\text{Cellule} \rightarrow \text{Tessuti} \rightarrow \text{Organi} \rightarrow \text{Sistemi} \rightarrow \text{Organismo}
\]
Nonostante la diversità morfologica e funzionale, le cellule condividono componenti strutturali di base, adattate alle specifiche esigenze fisiologiche.

\section{Procarioti ed eucarioti}
\begin{definizione}[Cellula eucariotica]
Una cellula \textbf{eucariotica} è caratterizzata dalla presenza di un nucleo delimitato da membrana, contenente il materiale genetico organizzato in cromosomi lineari, e di organelli membranosi specializzati.
\end{definizione}
La distinzione fondamentale tra i due domini cellulari è sintetizzata nella Tabella~\ref{tab:procarioti-eucarioti}.
\begin{table}[h!]
\centering
\caption{Confronto tra cellule procariotiche ed eucariotiche}
\label{tab:procarioti-eucarioti}
\begin{tabular}{p{0.45\textwidth} p{0.45\textwidth}}
\hline
\textbf{Procarioti} & \textbf{Eucarioti} \\
\hline
DNA circolare libero nel citoplasma & DNA lineare racchiuso nel nucleo \\
Assenza di organelli membranosi & Presenza di organelli membranosi \\
Dimensioni tipiche: 1--10 µm & Dimensioni tipiche: 10--100 µm \\
Organismi: batteri & Organismi: animali, piante, funghi \\
\hline
\end{tabular}
\end{table}

\section{Composizione chimica della cellula}
La maggior parte delle cellule presenta un contenuto idrico superiore al 70\%. Il \textbf{peso secco} (privato dell'acqua) è costituito principalmente da macromolecole biologiche, come riportato in Tabella~\ref{tab:composizione-cellula}.
\begin{table}[h!]
\centering
\caption{Composizione percentuale del peso secco cellulare}
\label{tab:composizione-cellula}
\begin{tabular}{lc}
\hline
\textbf{Componente} & \textbf{Peso secco (\%)} \\
\hline
Proteine & 75 \\
Lipidi & 12,5 \\
Glucidi & 5 \\
Acidi nucleici (DNA, RNA) & 5 \\
Composti inorganici / sali minerali & 2,5 \\
\hline
\end{tabular}
\end{table}
Le biomolecole si organizzano in \textbf{macromolecole} e successivamente in \textbf{complessi sopramolecolari}, che costituiscono le strutture funzionali della cellula.

\section{La membrana plasmatica}
La \textbf{membrana plasmatica} è una struttura dinamica che svolge funzioni essenziali:
\begin{itemize}
    \item \textbf{Separazione}: delimita il citoplasma dall'ambiente extracellulare;
    \item \textbf{Filtro selettivo}: regola il passaggio di sostanze in entrata e in uscita;
    \item \textbf{Comunicazione}: ospita recettori per segnali ormonali e molecole di adesione;
    \item \textbf{Interazione}: media i contatti fisici con altre cellule e con la matrice extracellulare.
\end{itemize}
\subsection{Struttura e composizione}
La membrana è costituita principalmente da un \textbf{doppio strato fosfolipidico}, in cui le code idrofobiche sono orientate verso l'interno e le teste polari verso l'ambiente acquoso. A questo si associano:
\begin{itemize}
    \item \textbf{Colesterolo}: modula la fluidità di membrana in risposta alle variazioni termiche;
    \item \textbf{Glicolipidi}: coinvolti nel riconoscimento cellulare;
    \item \textbf{Proteine di membrana}, distinte in:
    \begin{itemize}
        \item \textit{Integrali} (o intrinseche): attraversano completamente il doppio strato;
        \item \textit{Periferiche} (o estrinseche): associate a una delle due facce della membrana.
    \end{itemize}
\end{itemize}

\section{Trasporto attraverso le membrane}
Il passaggio di molecole attraverso la membrana avviene tramite meccanismi distinti, classificabili in base alla direzione del gradiente e al consumo energetico.
\subsection{Trasporto passivo}
Nel \textbf{trasporto passivo}, le molecole si muovono \textit{secondo gradiente} di concentrazione (da zone a concentrazione maggiore a zone a concentrazione minore), senza dispendio energetico.
\begin{itemize}
    \item \textbf{Diffusione semplice}: molecole apolari o gas (\(\mathrm{O_2}\), \(\mathrm{CO_2}\), \(\mathrm{N_2}\)) attraversano direttamente il doppio strato lipidico;
    \item \textbf{Diffusione facilitata}: molecole polari, cariche o di grandi dimensioni (es. glucosio) richiedono proteine di trasporto specifiche (es. trasportatori GLUT).
\end{itemize}
\subsection{Trasporto attivo}
Nel \textbf{trasporto attivo}, le molecole vengono spostate \textit{contro gradiente} di concentrazione, richiedendo energia (generalmente sotto forma di ATP) e proteine carrier specializzate.
\begin{definizione}[Tipologie di trasporto mediato]
\begin{itemize}
    \item \textbf{Uniporto}: trasporto di una singola specie molecolare;
    \item \textbf{Simporto}: trasporto contemporaneo di due specie nella stessa direzione;
    \item \textbf{Antiporto}: trasporto contemporaneo di due specie in direzioni opposte.
\end{itemize}
\end{definizione}
\subsection{Proteine recettoriali}
Alcune proteine di membrana fungono da \textbf{recettori}, capaci di legare specifici ligandi extracellulari (es. ormoni) e trasdurre il segnale all'interno della cellula, attivando cascate di risposta metabolica.
\begin{esempio}[Segnalazione insulinica]
L'insulina, secreta dal pancreas, si lega al proprio recettore di membrana sulle cellule muscolari, innescando una cascata di fosforilazioni che promuove l'ingresso di glucosio e la sintesi di glicogeno.
\end{esempio}

\section{Organelli cellulari}
Oltre alla membrana plasmatica, le cellule eucariotiche presentano un sistema di membrane interne che delimitano \textbf{organelli} specializzati, ciascuno con composizione lipidica e proteica specifica.
\subsection{Il nucleo}
\begin{itemize}
    \item Contiene il materiale genetico (DNA lineare) organizzato in \textbf{cromosomi}, complessati con proteine basiche dette \textbf{istoni};
    \item È delimitato da un \textbf{involucro nucleare} a doppia membrana, perforato da \textbf{pori nucleari} che regolano il traffico nucleo-citoplasmatico;
    \item Le cellule umane somatiche possiedono 46 cromosomi (23 coppie);
    \item Eccezioni: i globuli rossi maturi sono anucleati; le fibre muscolari scheletriche sono multinucleate.
\end{itemize}
\subsection{Il citoplasma e il citoscheletro}
Il \textbf{citoplasma} comprende tutto il contenuto cellulare escluso il nucleo ed è costituito da:
\begin{itemize}
    \item \textbf{Citosol}: fase acquosa contenente ioni, metaboliti e macromolecole;
    \item \textbf{Organelli}: strutture membranose specializzate;
    \item \textbf{Citoscheletro}: rete dinamica di filamenti proteici che conferisce forma, stabilità e capacità di movimento alla cellula.
\end{itemize}
\begin{table}[h!]
\centering
\caption{Componenti del citoscheletro}
\label{tab:citoscheletro}
\begin{tabular}{p{0.25\textwidth} p{0.35\textwidth} p{0.35\textwidth}}
\hline
\textbf{Tipo} & \textbf{Composizione} & \textbf{Funzioni principali} \\
\hline
Microfilamenti & Actina & Supporto meccanico, contrazione muscolare, motilità cellulare \\
Filamenti intermedi & Proteine fibrose (es. cheratine, vimentina) & Ancoraggio di organelli, resistenza meccanica \\
Microtubuli & Tubulina & Determinazione della forma, trasporto intracellulare, fuso mitotico \\
\hline
\end{tabular}
\end{table}
\subsection{Reticolo endoplasmatico}
Il \textbf{reticolo endoplasmatico} (RE) è una rete di membrane continue con l'involucro nucleare, distinta in due domini funzionali:
\begin{itemize}
    \item \textbf{RE rugoso}: presenta ribosomi associati alla superficie citosolica; è sede di sintesi e modificazione post-traduzionale delle proteine destinate a secrezione o a organelli membranosi;
    \item \textbf{RE liscio}: privo di ribosomi; coinvolto nella biosintesi lipidica, nel metabolismo dei carboidrati e nella detossificazione. Nel muscolo scheletrico, il RE liscio specializzato (\textbf{reticolo sarcoplasmatico}) regola il rilascio di ioni \(\mathrm{Ca^{2+}}\) necessari per la contrazione.
\end{itemize}
\subsection{Apparato di Golgi}
L'\textbf{apparato di Golgi} è costituito da una serie di cisterne membranose impilate, organizzate in due polarità:
\begin{itemize}
    \item \textbf{Lato \textit{cis}}: riceve vescicole di trasporto dal RE rugoso;
    \item \textbf{Lato \textit{trans}}: rilascia vescicole mature verso la membrana plasmatica, i lisosomi o altri compartimenti.
\end{itemize}
Funzioni principali: smistamento e indirizzamento delle proteine neosintetizzate, modificazioni post-traduzionali (glicosilazione, solfatazione), sintesi di carboidrati complessi e glicolipidi.
\subsection{Mitocondri}
I \textbf{mitocondri} sono gli organelli deputati alla produzione di energia cellulare sotto forma di ATP, tramite la fosforilazione ossidativa. Struttura a doppia membrana (esterna permeabile, interna ricca di creste e cardiolipina), contenente DNA mitocondriale circolare e ribosomi propri. Ereditati per via materna.
\subsection{Lisosomi}
I \textbf{lisosomi} sono vescicole membranose contenenti \textbf{enzimi idrolitici} attivi a pH acido (\(\mathrm{pH} \approx 4.5\text{--}5.0\)), deputati alla degradazione di macromolecole endocitate o di componenti cellulari danneggiati (autofagia). La membrana è protetta da un rivestimento glicoproteico interno.

% =========================================================
% Fine Capitolo 2
% =========================================================

% =========================================================
% CAPITOLO 3: Amminoacidi
% =========================================================
\chapter{Amminoacidi}
\section{Struttura generale e proprietà chimiche}
Gli \textbf{amminoacidi} sono molecole organiche caratterizzate dalla presenza di almeno due gruppi funzionali specifici: un \textbf{gruppo carbossilico} (\(-\mathrm{COOH}\)) e un \textbf{gruppo amminico} (\(-\mathrm{NH_2}\)). Negli amminoacidi costituenti le proteine, entrambi i gruppi sono legati allo stesso atomo di carbonio, denominato \(\alpha\)-carbonio (\(\alpha\)-C). Per questa ragione vengono definiti \(\alpha\)-amminoacidi.
Oltre al gruppo carbossilico e amminico, il \(\alpha\)-C è sempre legato a un atomo di idrogeno (\(\mathrm{H}\)) e a un gruppo variabile indicato con \(\mathrm{R}\), detto \textbf{catena laterale}. La composizione, la dimensione e la carica del gruppo \(\mathrm{R}\) differiscono per ciascun amminoacido e ne determinano le proprietà fisico-chimiche distintive, come la solubilità in acqua o la reattività.

\section{Chiralità e stereoisomeria}
Con l'eccezione della \textbf{glicina}, il cui gruppo \(\mathrm{R}\) è costituito da un singolo atomo di idrogeno, il \(\alpha\)-C di tutti gli amminoacidi è legato a quattro gruppi differenti. Di conseguenza, il \(\alpha\)-C rappresenta un \textbf{centro chirale} e i quattro sostituenti possono disporsi nello spazio in due configurazioni differenti, generando due stereoisomeri noti come \textbf{enantiomeri}. Gli enantiomeri sono immagini speculari non sovrapponibili l'una dell'altra.
Le due configurazioni sono designate con la notazione \textbf{D} (destrogiro) e \textbf{L} (levogiro), convenzionalmente definite in riferimento alla configurazione della gliceraldeide. Nei sistemi biologici, \textbf{tutti gli amminoacidi incorporati nelle proteine appartengono alla serie L}. Gli amminoacidi della serie D sono rari e si ritrovano esclusivamente in alcuni peptidi batterici o in contesti metabolici specifici. Questa omogeneità stereochimica è dovuta alla specificità degli enzimi biosintetici, che catalizzano reazioni stereospecifiche producendo esclusivamente L-amminoacidi.

\section{Ionizzazione e forma zwitterionica}
I gruppi carbossilico e amminico, nonché alcune catene laterali, sono \textbf{ionizzabili} e il loro stato di protonazione dipende dal pH dell'ambiente:
\begin{itemize}
    \item \textbf{A pH acido}: entrambi i gruppi sono protonati (\(-\mathrm{COOH}\) e \(-\mathrm{NH_3^+}\)); la molecola presenta una carica netta positiva.
    \item \textbf{A pH fisiologico} (\(\approx 7.4\)): il gruppo carbossilico perde un protone (\(-\mathrm{COO^-}\)), mentre il gruppo amminico rimane protonato (\(-\mathrm{NH_3^+}\)). La molecola possiede simultaneamente una carica positiva e una negativa, risultando elettricamente neutra. Questa forma è detta \textbf{zwitterione}. Fanno eccezione gli amminoacidi con catene laterali ionizzabili (es. glutammato, lisina).
    \item \textbf{A pH basico}: anche il gruppo amminico perde un protone (\(-\mathrm{NH_2}\)); la molecola assume una carica netta negativa.
\end{itemize}

\section{Classificazione in base alla catena laterale}
Le proteine sono costituite da 20 amminoacidi standard (più la selenocisteina, considerata il 21°), identificabili tramite abbreviazioni di tre lettere o simboli di una lettera. Sulla base delle proprietà chimiche del gruppo \(\mathrm{R}\), vengono classificati in cinque categorie principali.

\subsection{Amminoacidi non polari (idrofobici)}
Presentano catene laterali apolari, insolubili in acqua, e tendono a localizzarsi all'interno delle proteine globulari per minimizzare il contatto con il solvente acquoso, stabilizzando la struttura tramite \textit{interazioni idrofobiche}.
\begin{itemize}
    \item \textbf{Glicina (Gly, G)}: il più semplice (\(\mathrm{R} = \mathrm{H}\)). Privo di chiralità, meno idrofobico del gruppo, non partecipa significativamente alle interazioni idrofobiche.
    \item \textbf{Alanina (Ala, A), Valina (Val, V), Leucina (Leu, L), Isoleucina (Ile, I)}: catene completamente idrocarburiche. Val, Leu e Ile sono noti come \textbf{BCAA} (amminoacidi a catena ramificata). Sono essenziali e possono essere ossidati dal muscolo scheletrico a scopo energetico. Leucina e isoleucina sono isomeri strutturali.
    \item \textbf{Metionina (Met, M)}: contiene zolfo legato tramite un ponte tioetere (\(-\mathrm{S-CH_3}\)), non polare.
    \item \textbf{Prolina (Pro, P)}: presenta un gruppo amminico secondario (imminico) che forma un anello ciclico rigido con il \(\alpha\)-C. Riduce la flessibilità conformazionale delle catene proteiche in cui è inserita.
\end{itemize}

\subsection{Amminoacidi aromatici}
Contengono anelli aromatici nelle catene laterali, conferendo caratteristiche generalmente non polari.
\begin{itemize}
    \item \textbf{Fenilalanina (Phe, F)}: completamente non polare.
    \item \textbf{Tirosina (Tyr, Y)}: possiede un gruppo ossidrilico (\(-\mathrm{OH}\)) fenolico che può formare legami idrogeno e subire \textbf{fosforilazione} (modificazione post-traduzionale regolatoria).
    \item \textbf{Triptofano (Trp, W)}: contiene un gruppo amminico secondario nell'anello indolico, capace di formare legami idrogeno.
\end{itemize}

\subsection{Amminoacidi polari non carichi}
Le catene laterali contengono gruppi funzionali capaci di formare legami idrogeno con l'acqua, rendendoli solubili. A pH fisiologico non presentano carica netta.
\begin{itemize}
    \item \textbf{Serina (Ser, S) e Treonina (Thr, T)}: polarità dovuta al gruppo ossidrilico (\(-\mathrm{OH}\)). La treonina differisce per la presenza di un gruppo metile. I gruppi \(-\mathrm{OH}\) di Ser, Thr e Tyr sono siti privilegiati per la \textbf{fosforilazione}, meccanismo chiave nella regolazione enzimatica e nella trasduzione del segnale.
    \item \textbf{Cisteina (Cys, C)}: contiene un gruppo sulfidrilico (\(-\mathrm{SH}\)) di modesta polarità. Due cisteine possono ossidarsi formando un \textbf{ponte disolfuro} (\(-\mathrm{S-S-}\)), legame covalente fondamentale per la stabilizzazione della struttura terziaria e quaternaria delle proteine.
    \item \textbf{Asparagina (Asn, N) e Glutammina (Gln, Q)}: contengono gruppi ammidici (\(-\mathrm{CONH_2}\)) che partecipano a reti di legami idrogeno. Sono le ammidi corrispondenti di acido aspartico e acido glutammico.
\end{itemize}

\subsection{Amminoacidi acidi (carichi negativamente)}
Presentano un gruppo carbossilico aggiuntivo nella catena laterale, ionizzabile a pH fisiologico. Sono pertanto carichi negativamente e fortemente idrofilici.
\begin{itemize}
    \item \textbf{Acido aspartico / Aspartato (Asp, D)}
    \item \textbf{Acido glutammico / Glutammato (Glu, E)}
\end{itemize}
Oltre al ruolo strutturale nelle proteine, questi amminoacidi svolgono funzioni metaboliche cruciali e partecipano alla trasduzione del segnale nervoso.

\subsection{Amminoacidi basici (carichi positivamente)}
Contengono gruppi basici che rimangono protonati a pH 7, conferendo carica positiva e massima idrofilicità.
\begin{itemize}
    \item \textbf{Lisina (Lys, K)}: catena lineare con gruppo amminico terminale.
    \item \textbf{Arginina (Arg, R)}: contiene il gruppo \textbf{guanidinico}, fortemente basico.
    \item \textbf{Istidina (His, H)}: possiede un anello imidazolico. Il suo pKa (\(\approx 6.0\)) lo rende parzialmente protonato a pH fisiologico, fungendo da tampone intracellulare e partecipando attivamente ai siti catalitici enzimatici.
\end{itemize}

\section{Classificazione nutrizionale e destino metabolico}
\subsection{Essenziali e non essenziali}
In base alla capacità biosintetica dell'organismo:
\begin{itemize}
    \item \textbf{Non essenziali}: sintetizzati \textit{de novo} dall'organismo.
    \item \textbf{Essenziali}: non sintetizzabili, devono essere introdotti con la dieta. I BCAA rientrano in questa categoria e sono tra i più abbondanti nelle proteine muscolari.
    \item \textbf{Condizionalmente essenziali}: in condizioni fisiologiche estreme (sviluppo, stress, patologie), la richiesta supera la capacità sintetica, rendendoli di fatto indispensabili dalla dieta.
\end{itemize}

\subsection{Destino metabolico}
In carenza di carboidrati e lipidi, gli amminoacidi possono essere catabolizzati a scopo energetico:
\begin{itemize}
    \item \textbf{Glucogenici}: la degradazione produce intermedi utilizzabili per la gluconeogenesi (es. Asp, Glu, Asn, Gln, His, Pro, Arg, Gly, Ala, Ser, Cys, Met, Val).
    \item \textbf{Chetogenici}: la degradazione produce acetil-CoA, precursore della chetogenesi (Leu, Lys).
    \item \textbf{Entrambi}: Phe, Tyr, Trp, Ile, Thr.
\end{itemize}
Il metabolismo degli amminoacidi (tranne i BCAA) avviene prevalentemente nel fegato, che li trasforma in substrati utilizzabili da altri tessuti. I BCAA sono invece ossidati direttamente nel muscolo scheletrico.

\section{Amminoacidi speciali e non proteici}
Oltre ai 20-21 amminoacidi proteinogenici, esistono derivati e modificazioni post-traduzionali con ruoli fisiologici distinti:
\begin{itemize}
    \item \textbf{Selenocisteina (Sec, U)}: 21° amminoacido proteinogenico. Analogo della cisteina con selenio al posto dello zolfo. Presente nelle \textbf{selenoproteine} (es. glutatione perossidasi) con funzioni antiossidanti.
    \item \textbf{Idrossilisina e Idrossiprolina}: derivano da lisina e prolina tramite idrossilazione post-traduzionale. Abbondanti nel collagene e nell'elastina; il gruppo \(-\mathrm{OH}\) favorisce la formazione di legami idrogeno e la glicosilazione.
    \item \textbf{L-Carnitina}: sintetizzata da lisina e metionina. Trasporta acidi grassi a catena lunga nella matrice mitocondriale per la \(\beta\)-ossidazione. La carenza provoca miopatie e accumulo lipidico.
    \item \textbf{GABA} (\(\gamma\)-aminobutirrico): derivato del glutammato, principale neurotrasmettitore inibitorio del SNC. Modula ansia, sonno e stimola il rilascio di GH.
    \item \textbf{Citrullina e Ornitina}: intermedi del ciclo dell'urea. Derivano dall'arginina; la citrullina è sottoprodotto della sintesi di ossido nitrico (NO), potente vasodilatatore. Utilizzate in ambito sportivo per migliorare il flusso ematico e ridurre la fatica.
    \item \textbf{\(\beta\)-alanina}: isomero dell'alanina con gruppo amminico in posizione \(\beta\). Precursore della carnosina (dipeptide con istidina), importante sistema tampone intramuscolare contro l'acidosi da esercizio. Componente del Coenzima A.
    \item \textbf{Taurina}: amminoacido solfonico (\(-\mathrm{SO_3H}\)). Abbondante nel muscolo, componente dei sali biliari e regolatore del volume cellulare. Presente negli energy drink, sebbene i benefici ergogenici siano dibattuti.
\end{itemize}

% =========================================================
% Fine Capitolo 3
% =========================================================

% =========================================================
% CAPITOLO 4: Le Proteine
% =========================================================
\chapter{Le Proteine}
\section{Introduzione e legame peptidico}
Le \textbf{proteine} sono macromolecole biologiche di dimensioni altamente variabili, comprese tra pochi kilodalton (kDa) e alcuni megadalton (MDa). Sono mediamente più estese di lipidi e carboidrati e le loro unità monomeriche sono i 20 amminoacidi proteinogenici.
Gli amminoacidi all'interno di una proteina sono uniti tramite il \textbf{legame peptidico}, un legame covalente che si forma tra il gruppo carbossilico (\(-\mathrm{COOH}\)) di un amminoacido e il gruppo amminico (\(-\mathrm{NH_2}\)) del successivo. La reazione libera una molecola di acqua e rappresenta quindi una \textbf{reazione di condensazione}. Il legame peptidico non coinvolge le catene laterali (\(\mathrm{R}\)).
Due amminoacidi uniti formano un \textbf{dipeptide}; tre formano un \textbf{tripeptide}. Per catene più lunghe si utilizzano le seguenti convenzioni:
\begin{itemize}
    \item \(< 20\) amminoacidi: \textbf{oligopeptide}
    \item \(20\text{--}100\) amminoacidi: \textbf{polipeptide}
    \item \(> 100\) amminoacidi: \textbf{proteina}
\end{itemize}
Nella rappresentazione lineare della sequenza, si adotta la convenzione di scrivere il gruppo amminico libero a sinistra e il gruppo carbossilico libero a destra. Il primo residuo è definito \textbf{N-terminale}, l'ultimo \textbf{C-terminale}. Poiché solo le estremità mantengono i gruppi funzionali liberi, le proprietà idrofiliche/idrofobiche della molecola dipendono principalmente dalla natura delle catene laterali.

\section{Struttura delle proteine}
La conformazione proteica viene descritta a quattro livelli gerarchici.

\subsection{Struttura primaria}
Rappresenta la sequenza lineare di amminoacidi, determinata univocamente dall'informazione genetica contenuta nel DNA. Il numero e l'ordine dei residui definiscono la funzione della proteina; anche una singola sostituzione (mutazione) può comprometterne l'attività o renderla patologica. La sequenza primaria contiene le informazioni necessarie per il ripiegamento tridimensionale, sebbene la previsione computazionale della struttura finale rimanga complessa.

\subsection{Struttura secondaria}
Corrisponde ai ripiegamenti locali ricorrenti della catena polipeptidica, stabilizzati da \textbf{legami idrogeno} regolari tra i gruppi ammidici e carbonilici dello scheletro peptidico. Le due conformazioni più comuni sono:
\begin{itemize}
    \item \textbf{\(\alpha\)-elica}: struttura a spirale stabilizzata da legami idrogeno intra-catena.
    \item \textbf{Foglietto \(\beta\) ripiegato}: disposizione planare di segmenti polipeptidici adiacenti, collegati da legami idrogeno inter-catena o intra-catena.
\end{itemize}
Combinazioni di \(\alpha\)-eliche e foglietti \(\beta\) formano \textbf{motivi strutturali} (strutture supersecondarie).

\subsection{Struttura terziaria}
È la conformazione tridimensionale complessiva della singola catena polipeptidica, ovvero la disposizione spaziale di tutti gli atomi. A differenza della struttura secondaria, la stabilità terziaria è garantita da interazioni tra catene laterali:
\begin{itemize}
    \item \textbf{Ponti disolfuro} (\(-\mathrm{S-S-}\)): legami covalenti tra residui di cisteina, molto stabili.
    \item \textbf{Interazioni idrofobiche}: i residui apolari si ripiegano verso l'interno della proteina per evitare il contatto con l'acqua, mentre i residui polari si dispongono sulla superficie.
    \item \textbf{Legami idrogeno} e \textbf{interazioni ioniche} tra gruppi carichi.
\end{itemize}
Una proteina funzionale assume la sua \textbf{conformazione nativa}; la perdita di tale struttura è definita \textbf{denaturazione}.
Le catene polipeptidiche lunghe (\(>200\) residui) si organizzano spesso in \textbf{domini strutturali}: unità globulari compatte, indipendenti dal punto di vista strutturale e funzionale, collegate da regioni flessibili.

\subsection{Struttura quaternaria}
Presente nelle proteine multimeriche, descrive l'associazione di due o più catene polipeptidiche (\textbf{subunità}). Le subunità possono essere identiche (\textbf{omomultimeriche}, es. omodimero) o diverse (\textbf{eteromultimeriche}, es. eterodimero). Le interazioni tra subunità sono dello stesso tipo di quelle che stabilizzano la struttura terziaria. Esempi noti sono l'\textbf{emoglobina} (tetramero: \(2\alpha + 2\beta\)) e la \textbf{miosina} (esamero: 2 catene pesanti + 4 leggere).

\section{Classificazione e proteine coniugate}
\subsection{Forma e solubilità}
\begin{itemize}
    \item \textbf{Proteine fibrose}: una dimensione prevale sulle altre. Presentano spesso un unico tipo di struttura secondaria ripetuta, sono insolubili in acqua e svolgono ruoli strutturali o di supporto (es. cheratine, collageno, elastina, seta).
    \item \textbf{Proteine globulari}: le tre dimensioni sono comparabili, assumendo forma compatta. Sono generalmente solubili e svolgono funzioni dinamiche (enzimi, trasportatori, ormoni, anticorpi).
\end{itemize}

\subsection{Proteine coniugate}
Molte proteine contengono \textbf{gruppi prostetici} non amminoacidici, essenziali per la funzione. Possono essere legati covalentemente (es. carboidrati nelle glicoproteine, fosfati nelle fosfoproteine), associati strettamente in tasche idrofobiche (es. gruppo \textbf{eme} nell'emoglobina) o legati tramite interazioni ioniche (metalli nelle metalloproteine).

\section{Modifiche post-traduzionali}
Dopo la sintesi, le proteine possono subire modificazioni chimiche che ne regolano attività, localizzazione o stabilità:
\begin{itemize}
    \item \textbf{Fosforilazione}: aggiunta di un gruppo fosfato (da ATP) a residui di Serina, Treonina o Tirosina. Modifica conformazione e attività; è un meccanismo reversibile di regolazione enzimatica (es. attivazione della glicogeno fosforilasi in risposta al glucagone).
    \item \textbf{Acetilazione}: aggiunta di un gruppo acetile (\(-\mathrm{COCH_3}\)) a residui di Lisina.
    \item \textbf{Glicosilazione}: aggiunta di catene oligosaccaridiche. La N-glicosilazione avviene sull'Asparagina, la O-glicosilazione su Serina/Treonina. Cruciale per il targeting e la funzione delle proteine di membrana.
\end{itemize}

\section{Sintesi, turnover e degradazione}
\subsection{Sintesi proteica}
L'informazione fluisce dal DNA all'mRNA (\textbf{trascrizione}), che viene esportato nel citosol e tradotto dai ribosomi in catena polipeptidica (\textbf{traduzione}). Il codice genetico è letto in triplette (\textbf{codoni}); 61 codoni specificano amminoacidi (codice degenerato), 3 sono codoni di stop. La traduzione inizia generalmente dal codone AUG (Metionina). Il folding avviene spontaneamente o con l'assistenza di \textbf{chaperoni}. Il processo è energeticamente costoso e regolato finemente, ad esempio dalla proteina \textbf{mTOR}, attivata in condizioni di abbondanza nutrizionale e promotrice dell'ipertrofia muscolare.

\subsection{Turnover e degradazione}
Le proteine sono soggette a continuo \textbf{turnover} (sintesi + degradazione), che consuma fino al 20\% del metabolismo basale. Ogni proteina ha un'\textbf{emivita} variabile (minuti per enzimi epatici, settimane per proteine muscolari). La degradazione avviene principalmente tramite:
\begin{itemize}
    \item \textbf{Lisosomi}: degradano proteine extracellulari (endocitosi) o intracellulari/organelli (autofagia) in condizioni di digiuno o stress. Processo generalmente non selettivo, ottimale a pH acido.
    \item \textbf{Proteasoma}: complesso citosolico che degrada proteine intracellulari marcate con \textbf{ubiquitina} (ubiquitinazione). Processo ATP-dipendente e selettivo.
    \item \textbf{Digestione proteica}: scissione delle proteine alimentari in amminoacidi assorbibili tramite proteasi gastrointestinali.
\end{itemize}

\section{Funzioni biologiche}
Le proteine svolgono una vasta gamma di ruoli fisiologici:
\begin{enumerate}
    \item \textbf{Strutturale e contrattile}: Actina e miosina (contrazione muscolare), collageno e cheratina (supporto tissutale), distrofina (ancoraggio miofibrille al sarcolemma).
    \item \textbf{Trasporto e deposito}: Proteine di membrana (canali, trasportatori GLUT), emoglobina (trasporto \(\mathrm{O_2}\) tramite gruppo eme).
    \item \textbf{Ormonale}: Insulina, glucagone, GH (regolazione metabolica e crescita).
    \item \textbf{Segnalazione}: Recettori di membrana e cascate di chinasi intracellulari che trasducono stimoli esterni in risposte cellulari.
    \item \textbf{Immunitaria}: Immunoglobuline (anticorpi) che riconoscono e neutralizzano antigeni.
    \item \textbf{Regolazione}: Fattori di trascrizione (es. PGC-1\(\alpha\)) che modulano l'espressione genica in risposta a stimoli come l'esercizio.
    \item \textbf{Catalitica}: Enzimi (es. esochinasi) che accelerano le reazioni metaboliche abbassando l'energia di attivazione.
\end{enumerate}

\section{Stabilità e denaturazione}
La funzione proteica è strettamente legata alla conformazione nativa. Fattori come variazioni di pH, aumento di temperatura o stress ossidativo possono rompere i legami deboli che stabilizzano la struttura, portando a \textbf{denaturazione} (spesso irreversibile, come nella cottura dell'uovo). Una proteina denaturata perde la sua funzione e può risultare tossica, venendo quindi degradata.
Durante l'esercizio intenso, l'ambiente intracellulare diventa acido (\(\mathrm{pH} \approx 6.5\)) e ipertermico (\(\approx 40^\circ\mathrm{C}\)). Per contrastare questi effetti, le cellule esprimono \textbf{Heat Shock Proteins (HSPs)}: chaperoni molecolari che riconoscono le proteine misfolded, ne facilitano il refolding e promuovono una risposta adattativa che aumenta la resilienza cellulare a stress futuri.

% =========================================================
% Fine Capitolo 4
% =========================================================

% =========================================================
% CAPITOLO 5: Trasporto e utilizzo dell'ossigeno
% =========================================================
\chapter{Trasporto e utilizzo dell'ossigeno}
\section{Introduzione: emoglobina e mioglobina}
L'\textbf{ossigeno} (\(\mathrm{O_2}\)) è l'elemento più abbondante nel nostro organismo, principalmente incorporato nell'acqua. Costituisce inoltre il 21\% dell'atmosfera terrestre. L'\(\mathrm{O_2}\) inspirato nei polmoni deve essere trasportato ai tessuti per sostenere il metabolismo ossidativo mitocondriale.
Nel corpo umano, due proteine specializzate legano reversibilmente l'ossigeno:
\begin{itemize}
    \item \textbf{Emoglobina (Hb)}: proteina tetramerica presente negli eritrociti (globuli rossi), deputata al trasporto di \(\mathrm{O_2}\) dai polmoni ai tessuti e di \(\mathrm{CO_2}\) dai tessuti ai polmoni.
    \item \textbf{Mioglobina (Mb)}: proteina monomerica localizzata nelle cellule muscolari, funge da riserva di \(\mathrm{O_2}\) rilasciandolo in condizioni di ipossia per mantenere attivo il metabolismo aerobico.
\end{itemize}

\section{Struttura delle emoproteine}
Sia l'emoglobina che la mioglobina sono \textbf{emoproteine}, ovvero contengono un gruppo prostetico \textbf{eme} legato non covalentemente alla catena polipeptidica. L'eme contiene un atomo di ferro nello stato di ossidazione ferroso (\(\mathrm{Fe^{2+}}\)), responsabile del legame reversibile con l'\(\mathrm{O_2}\).

\subsection{Il gruppo eme}
L'eme è una protoporfirina IX, struttura planare derivante dalla porfirina, caratterizzata da quattro anelli pirrolici con atomi di azoto all'apice. Al centro della protoporfirina è coordinato un singolo atomo di \(\mathrm{Fe^{2+}}\).
I legami di coordinazione del ferro sono sei:
\begin{itemize}
    \item Quattro legami giacciono sul piano dell'anello porfirinico e coinvolgono gli atomi di azoto dei pirroli;
    \item Il quinto legame è occupato dall'azoto dell'anello imidazolico di un residuo di \textbf{istidina prossimale};
    \item Il sesto legame, perpendicolare al piano, è disponibile per il legame con l'\(\mathrm{O_2}\).
\end{itemize}
Un secondo residuo di \textbf{istidina distale} stabilizza il legame \(\mathrm{Fe-O_2}\) tramite un legame idrogeno, favorendo il posizionamento planare del ferro rispetto all'eme. Per prevenire l'ossidazione del ferro a \(\mathrm{Fe^{3+}}\) (che impedirebbe il legame con \(\mathrm{O_2}\)), il gruppo eme è alloggiato in una tasca idrofobica della proteina.

\subsection{Differenze strutturali tra Hb e Mb}
\begin{table}[h!]
\centering
\caption{Confronto tra emoglobina e mioglobina}
\label{tab:hb-mb}
\begin{tabular}{p{0.45\textwidth} p{0.45\textwidth}}
\hline
\textbf{Emoglobina} & \textbf{Mioglobina} \\
\hline
Tetramero (\(2\alpha + 2\beta\) nell'adulto) & Monomero (singola catena) \\
4 gruppi eme, lega fino a 4 \(\mathrm{O_2}\) & 1 gruppo eme, lega 1 \(\mathrm{O_2}\) \\
Curva di saturazione sigmoide (cooperatività) & Curva di saturazione iperbolica \\
Bassa affinità a bassa \(p\mathrm{O_2}\) & Alta affinità anche a bassa \(p\mathrm{O_2}\) \\
Trasporto sistemico di \(\mathrm{O_2}\) & Riserva intramuscolare di \(\mathrm{O_2}\) \\
\hline
\end{tabular}
\end{table}

\section{Legame dell'ossigeno e cooperatività}
La mioglobina lega l'\(\mathrm{O_2}\) con alta affinità anche a pressioni parziali (\(p\mathrm{O_2}\)) relativamente basse, mostrando una curva di saturazione \textbf{iperbolica}. Ciò la rende ideale per immagazzinare \(\mathrm{O_2}\) nel muscolo.
L'emoglobina, invece, presenta una curva di saturazione \textbf{sigmoide}, riflettente il legame cooperativo tra le sue quattro subunità. Questo comportamento è spiegato dal modello degli stati conformazionali:
\begin{itemize}
    \item \textbf{Stato T (teso)}: conformazione a bassa affinità per \(\mathrm{O_2}\), predominante in condizioni di bassa \(p\mathrm{O_2}\) (tessuti);
    \item \textbf{Stato R (rilassato)}: conformazione ad alta affinità, favorita dal legame dell'\(\mathrm{O_2}\) e predominante nei polmoni.
\end{itemize}
Il legame dell'\(\mathrm{O_2}\) a una subunità induce un cambiamento conformazionale che facilita il legame alle subunità adiacenti (\textbf{cooperatività omotropica}). Nei polmoni (alta \(p\mathrm{O_2}\)), l'Hb si satura rapidamente; nei tessuti (bassa \(p\mathrm{O_2}\)), il rilascio è favorito dalla transizione allo stato T.

\section{Regolazione allosterica dell'affinità per l'ossigeno}
L'affinità dell'emoglobina per l'\(\mathrm{O_2}\) è modulata da fattori allosterici eterotropici che agiscono su siti distinti dal sito di legame dell'\(\mathrm{O_2}\):

\subsection{Effetto Bohr (pH)}
L'aumento della concentrazione di ioni \(\mathrm{H^+}\) (diminuzione del pH) riduce l'affinità dell'Hb per l'\(\mathrm{O_2}\), favorendone il rilascio nei tessuti metabolicamente attivi. Durante l'esercizio intenso, la produzione di acido lattico e di \(\mathrm{CO_2}\) abbassa il pH tissutale:
\[
\mathrm{CO_2 + H_2O \rightleftharpoons H_2CO_3 \rightleftharpoons HCO_3^- + H^+}
\]
Gli ioni \(\mathrm{H^+}\) protonano residui di istidina, stabilizzando lo stato T e promuovendo il rilascio di \(\mathrm{O_2}\). Nei polmoni, l'eliminazione di \(\mathrm{CO_2}\) e l'aumento del pH favoriscono il ricarico di \(\mathrm{O_2}\).

\subsection{Anidride carbonica (\(\mathrm{CO_2}\))}
Oltre all'effetto indiretto tramite il pH, la \(\mathrm{CO_2}\) lega direttamente l'N-terminale delle catene dell'Hb formando \textbf{carbammati}, stabilizzando lo stato T e riducendo l'affinità per l'\(\mathrm{O_2}\). Circa il 20\% della \(\mathrm{CO_2}\) prodotta dai tessuti è trasportata dall'Hb sotto questa forma.
\begin{nota}[Monossido di carbonio]
Il monossido di carbonio (\(\mathrm{CO}\)) lega il ferro dell'eme con affinità ~200 volte superiore all'\(\mathrm{O_2}\), formando carbossiemoglobina stabile e impedendo il trasporto di ossigeno (ipossia anemica).
\end{nota}

\subsection{Temperatura}
L'aumento della temperatura riduce l'affinità dell'Hb per l'\(\mathrm{O_2}\). Durante l'esercizio intenso, la temperatura muscolare può raggiungere i 40°C, favorendo il rilascio di \(\mathrm{O_2}\) a livello capillare per sostenere il metabolismo ossidativo.

\subsection{2,3-Bisfosfoglicerato (2,3-BPG)}
Il \textbf{2,3-BPG} è un metabolita glicolitico prodotto dagli eritrociti che lega esclusivamente la deossiemoglobina in una cavità centrale tra le quattro subunità, stabilizzando lo stato T e riducendo l'affinità per l'\(\mathrm{O_2}\).
\begin{itemize}
    \item In condizioni di ipossia (alta quota), la sintesi di 2,3-BPG aumenta, facilitando il rilascio tissutale di \(\mathrm{O_2}\) (adattamento a breve termine).
    \item L'emoglobina fetale (\(\mathrm{HbF} = \alpha_2\gamma_2\)) ha minore affinità per il 2,3-BPG rispetto all'Hb adulta, garantendo un'affinità maggiore per l'\(\mathrm{O_2}\) e permettendo il trasferimento di ossigeno dal sangue materno a quello fetale.
\end{itemize}

\section{Adattamenti all'ipossia e considerazioni cliniche}
\subsection{Adattamenti fisiologici}
Oltre all'aumento di 2,3-BPG, l'esposizione prolungata all'ipossia stimola la produzione renale di \textbf{eritropoietina (EPO)}, un ormone glicoproteico che promuove l'eritropoiesi nel midollo osseo, aumentando la capacità di trasporto di \(\mathrm{O_2}\) del sangue (adattamento a lungo termine).

\subsection{Beta-talassemia (anemia mediterranea)}
La \textbf{beta-talassemia} è una malattia genetica causata da mutazioni che riducono o aboliscono la sintesi delle catene \(\beta\) dell'emoglobina. Il risultato è la produzione di globuli rossi difettosi, anemia emolitica e necessità di trasfusioni frequenti. La comprensione della regolazione dell'Hb ha implicazioni terapeutiche per queste patologie.

\section{Sintesi funzionale}
\begin{itemize}
    \item Emoglobina e mioglobina sono emoproteine essenziali per il trasporto e lo stoccaggio dell'\(\mathrm{O_2}\).
    \item La struttura tetramerica dell'Hb e la cooperatività omotropica consentono un rilascio efficiente di \(\mathrm{O_2}\) nei tessuti.
    \item Fattori allosterici eterotropici (pH, \(\mathrm{CO_2}\), temperatura, 2,3-BPG) modulano finemente l'affinità dell'Hb per l'\(\mathrm{O_2}\) in risposta alle esigenze metaboliche.
    \item Adattamenti fisiologici (2,3-BPG, EPO) e patologie (talassemia) illustrano l'importanza clinica della regolazione del trasporto di ossigeno.
\end{itemize}

% =========================================================
% Fine Capitolo 5
% =========================================================

% =========================================================
% CAPITOLO 6: Enzimi
% =========================================================

\chapter{Enzimi}

\section{Introduzione e ruolo catalitico}
Gli \textbf{enzimi} sono proteine (con rare eccezioni rappresentate da RNA catalitici detti \textbf{ribozimi}) che svolgono la funzione di \textbf{catalizzatori biologici}. La loro presenza è indispensabile per accelerare le reazioni chimiche cellulari che, in condizioni fisiologiche, avverrebbero a velocità incompatibili con la vita.
Gli enzimi presentano tre caratteristiche fondamentali:
\begin{itemize}
    \item \textbf{Non vengono consumati} durante la reazione e possono quindi catalizzare numerosi cicli catalitici.
    \item \textbf{Non alterano l'equilibrio} della reazione: non influenzano la direzione termodinamica, ma accelerano il raggiungimento dell'equilibrio.
    \item \textbf{Aumentano la velocità} sia della reazione diretta che di quella inversa, fino a un fattore di \(10^{10}\) volte o superiore.
\end{itemize}
Un esempio classico è la reazione catalizzata dalla lattato deidrogenasi (LDH):
\[
\text{Piruvato} + \text{NADH} \rightleftharpoons \text{Lattato} + \text{NAD}^+
\]

\section{Energia di attivazione e meccanismo d'azione}
Una reazione esoergonica (\(\Delta G < 0\)) è termodinamicamente favorita, ma ciò non garantisce che avvenga rapidamente. Per procedere, i reagenti devono superare una barriera energetica, raggiungendo uno \textbf{stato di transizione} ad alta energia. L'energia necessaria per raggiungere questo stato è definita \textbf{energia di attivazione} (\(E_a\)).
È possibile immaginare l'\(E_a\) come l'energia necessaria per spingere un masso (reagenti) sulla sommità di una collina (stato di transizione), prima che rotoli a valle (prodotti). Più alta è la \(E_a\), minore è la probabilità che i reagenti si orientino casualmente in modo favorevole alla reazione, risultando in una cinetica estremamente lenta.
La funzione esclusiva degli enzimi è \textbf{abbassare l'energia di attivazione}, stabilizzando lo stato di transizione e permettendo alla reazione di procedere a velocità fisiologicamente rilevanti, senza alterare il \(\Delta G\) complessivo.
\begin{esempio}[Fosforilazione del glucosio]
La prima tappa della glicolisi (glucosio \(\rightarrow\) glucosio-6-fosfato) ha un \(\Delta G\) negativo, ma un'elevata \(E_a\). In assenza di catalisi, la reazione sarebbe trascurabile. L'enzima \textbf{esochinasi} abbassa la \(E_a\) favorendo l'interazione tra glucosio e ATP, consentendo la rapida formazione dei prodotti.
\end{esempio}

\section{Specificità e modelli di riconoscimento del substrato}
Per abbassare l'\(E_a\), gli enzimi legano reversibilmente i substrati formando un \textbf{complesso enzima-substrato (ES)}. Il legame avviene in una regione specifica della proteina detta \textbf{sito attivo}, composta da residui amminoacidici coinvolti sia nel riconoscimento del substrato sia nella catalisi. Durante il processo si formano intermedi transitori (ES ed EP) prima del rilascio del prodotto.

\subsection{Modelli di interazione}
Storicamente, il \textbf{modello chiave-serratura} ipotizzava una complementarità geometrica rigida e preformata tra enzima e substrato. Tuttavia, un legame eccessivamente stabile impedirebbe il successivo rilascio del prodotto.
Il modello attualmente accettato è quello dell'\textbf{adattamento indotto}: il sito attivo non è perfettamente complementare al substrato nello stato libero. L'interazione iniziale (mediata da legami deboli e transitori) induce un cambiamento conformazionale nell'enzima che ottimizza il posizionamento del substrato, favorisce la catalisi e agevola il successivo distacco del prodotto. Un esempio è la chiusura conformazionale dell'esochinasi attorno al glucosio dopo il legame.

\subsection{Stereospecificità}
Oltre alla specificità chimica, gli enzimi sono altamente \textbf{stereospecifici}. Poiché le proteine sono sintetizzate esclusivamente da L-amminoacidi, i siti attivi presentano un'architettura chirale asimmetrica. Di conseguenza, riconoscono e trasformano selettivamente uno specifico stereoisomero, e rilasciano prodotti con configurazione spaziale definita. Esempio: la lattato deidrogenasi muscolare è specifica per l'\textbf{L-lattato}.

\section{Cinetica enzimatica e fattori di modulazione}
L'attività enzimatica è finemente modulata dalla cellula per adattare il flusso metabolico alle richieste energetiche. I principali fattori che influenzano la velocità di reazione sono: concentrazione del substrato, pH, temperatura e concentrazione dell'enzima.

\subsection{Concentrazione del substrato e costante di Michaelis}
In una reazione non catalizzata, la velocità aumenta linearmente con la concentrazione del substrato \([S]\). Nelle reazioni enzimatiche, la velocità aumenta linearmente solo a basse concentrazioni; oltre una certa soglia, la curva diventa \textbf{iperbolica} fino a raggiungere un plateau, la \textbf{velocità massima} (\(V_{\max}\)), quando tutti i siti attivi sono saturi.
La relazione è descritta dalla \textbf{costante di Michaelis} (\(K_M\)), definita come la concentrazione di substrato alla quale la velocità di reazione è pari a \(V_{\max}/2\).
\begin{itemize}
    \item \(K_M\) \textbf{bassa} \(\rightarrow\) alta affinità enzima-substrato.
    \item \(K_M\) è specifica per ogni coppia enzima-substrato.
    \item In condizioni fisiologiche, \([S]\) è generalmente \(< K_M\), collocando l'enzima nella zona lineare della curva: piccole variazioni di substrato si traducono in ampie variazioni della velocità di reazione, garantendo una risposta metabolica sensibile.
\end{itemize}

\subsection{Effetto del pH e della temperatura}
\begin{itemize}
    \item \textbf{pH}: La maggior parte degli enzimi presenta un optimum attorno a pH 7. Eccezioni includono enzimi digestivi gastrici (pH 1--2) e lisosomiali (pH 4,5--5). Deviazioni dal pH ottimale alterano lo stato di ionizzazione dei residui del sito attivo, riducendo l'efficienza catalitica. Durante esercizio intenso, il pH muscolare può scendere a 6,5, inibendo enzimi metabolici e contribuendo all'insorgenza della fatica.
    \item \textbf{Temperatura}: L'aumento termico incrementa l'energia cinetica e la frequenza delle collisioni molecolari, accelerando la reazione fino a \(\approx 50^\circ\mathrm{C}\). Oltre questa soglia, i legami deboli che stabilizzano la struttura terziaria si rompono, causando \textbf{denaturazione} e perdita irreversibile dell'attività. Il riscaldamento pre-esercizio sfrutta questo principio per ottimizzare l'efficienza enzimatica muscolare.
\end{itemize}

\subsection{Concentrazione dell'enzima}
Quando l'enzima opera in condizioni di saturazione (\([S] \gg K_M\)), la velocità di reazione è direttamente proporzionale alla concentrazione dell'enzima \([E]\), mentre la \(K_M\) rimane invariata. La capacità delle cellule di aumentare l'espressione proteica enzimatica è il fondamento degli \textbf{adattamenti metabolici all'allenamento}, in particolare di tipo aerobico.

\section{Cofattori, coenzimi e vitamine}
Molti enzimi richiedono molecole non proteiche per esercitare la loro funzione catalitica:
\begin{itemize}
    \item \textbf{Cofattori}: ioni inorganici essenziali (es. \(\mathrm{Cu^{2+}}, \mathrm{K^+}, \mathrm{Fe^{2+}}, \mathrm{Mg^{2+}}, \mathrm{Mn^{2+}}, \mathrm{Zn^{2+}}\)).
    \item \textbf{Coenzimi}: molecole organiche che partecipano direttamente alla reazione, spesso subendo modificazioni temporanee (agendo come secondi substrati). Esempio: \(\mathrm{NAD^+}\) nella lattato deidrogenasi.
    \item \textbf{Gruppo prostetico}: cofattore o coenzima legato stabilmente (spesso covalentemente) all'enzima (es. gruppo eme nell'emoglobina).
\end{itemize}
Molte \textbf{vitamine} (specie del gruppo B) sono precursori essenziali di coenzimi. Carenze vitaminiche compromettono l'attività enzimatica e causano patologie metaboliche. I coenzimi \(\mathrm{NAD^+}\) e \(\mathrm{FAD}\) sono fondamentali come accettori/donatori di elettroni nelle reazioni redox mediate dalle deidrogenasi.

\section{Classificazione e nomenclatura enzimatica}
Gli enzimi sono classificati in \textbf{6 classi principali} in base al tipo di reazione catalizzata. La maggior parte presenta il suffisso \textit{-asi}. Ogni enzima è identificato da un codice EC (Enzyme Commission) a 4 numeri (Classe.Sottoclasse.Sottosottoclasse.Seriale).

\begin{table}[h!]
\centering
\caption{Classificazione principale degli enzimi}
\label{tab:enzimi-classi}
\begin{tabular}{p{0.25\textwidth} p{0.70\textwidth}}
\hline
\textbf{Classe} & \textbf{Tipo di reazione} \\
\hline
1. Ossidoreduttasi & Reazioni redox, trasferimento di atomi di idrogeno/elettroni \\
2. Transferasi & Trasferimento di gruppi funzionali (es. chinasi: trasferiscono fosfato) \\
3. Idrolasi & Rottura di legami per aggiunta di \(\mathrm{H_2O}\) (idrolisi) \\
4. Liasi & Eliminazione di gruppi per formare doppi legami, o addizione inversa \\
5. Isomerasi & Riarrangiamento intramolecolare (isomerizzazione) \\
6. Ligasi & Formazione di nuovi legami (C-C, C-O, C-N) accoppiata a idrolisi di ATP \\
\hline
\end{tabular}
\end{table}

\section{Regolazione e inibizione dell'attività enzimatica}
Oltre a pH e temperatura, l'attività enzimatica è regolata finemente attraverso due meccanismi principali: modificazioni covalenti e regolazione allosterica.

\subsection{Modificazioni covalenti}
Consistono nell'aggiunta o rimozione reversibile di gruppi funzionali che alterano la conformazione e l'attività dell'enzima. Il meccanismo più comune è la \textbf{fosforilazione}: l'aggiunta di un gruppo fosfato (donato dall'ATP) a residui di Serina, Treonina o Tirosina, catalizzata da \textbf{chinasi}. La rimozione è operata da \textbf{fosfatasi} (defosforilazione). Sebbene la fosforilazione spesso attivi l'enzima, esistono eccezioni (es. inibizione della glicogeno sintasi).

\subsection{Regolazione allosterica}
Molecole chiamate \textbf{effettori allosterici} si legano a siti regolatori distinti dal sito attivo, inducendo cambiamenti conformazionali o elettrostatici che modulano l'affinità per il substrato. Gli effettori possono essere \textbf{positivi} (attivatori) o \textbf{negativi} (inibitori). Questo meccanismo è cruciale per il controllo fine della produzione di energia e biomolecole.

\subsection{Inibizione enzimatica}
L'attività può essere compromessa da molecole inibitrici:
\begin{itemize}
    \item \textbf{Inibizione irreversibile}: legame covalente o estremamente stabile che inattiva permanentemente l'enzima (tipica di farmaci o tossine).
    \item \textbf{Inibizione reversibile}: legame non covalente transitorio, fondamentale per la regolazione fisiologica.
    \begin{itemize}
        \item \textit{Competitiva}: l'inibitore compete con il substrato per il sito attivo (strutturalmente simile), riducendo l'efficienza senza alterare \(V_{\max}\) (superabile aumentando \([S]\)).
        \item \textit{Non competitiva}: l'inibitore si lega a un sito alternativo, alterando la conformazione e riducendo \(V_{\max}\) senza influenzare \(K_M\).
    \end{itemize}
\end{itemize}

% =========================================================
% Fine Capitolo 6
% =========================================================

\end{document}